\documentclass[11pt]{article}

\usepackage[margin=1.1in]{geometry}
\usepackage{amsmath,amssymb,amsthm,mathtools}
\usepackage{enumitem}
\usepackage{hyperref}
\usepackage[nameinlink,noabbrev]{cleveref}

\hypersetup{colorlinks=true,linkcolor=blue,citecolor=blue,urlcolor=blue}

\newtheorem{theorem}{Theorem}[section]
\newtheorem{lemma}[theorem]{Lemma}
\newtheorem{corollary}[theorem]{Corollary}
\newtheorem{proposition}[theorem]{Proposition}
\theoremstyle{remark}
\newtheorem{remark}[theorem]{Remark}

\providecommand{\doi}[1]{\href{https://doi.org/#1}{doi:#1}}
\renewcommand{\Re}{\operatorname{Re}}
\renewcommand{\Im}{\operatorname{Im}}
\newcommand{\D}{\mathcal D}
\newcommand{\tr}{\operatorname{tr}}

\title{Two questions of Jerrum and Patel on zeros of the independence polynomial}
\author{Yitzchak Shmalo\\
Einstein Institute of Mathematics, The Hebrew University of Jerusalem\\
Givat Ram, Jerusalem, Israel}
\date{}

\begin{document}
\maketitle

\begin{abstract}
Jerrum and Patel asked whether the claw-free graphs of bounded degree, or of bounded clique number, have an open
multivariate zero-free region for the independence polynomial containing $[0,\infty)$, and whether the graphs of
maximum degree $\Delta$ with no induced subdivided claw $S_{t,t,t}$ have a zero-free sector. For the first question
the answer is positive. Jerrum and Patel consider the parabolic region
$R(K)=\{z\in\mathbb C:(\Im z)^2<\Re z/K+1/(4K^2)\}$; we show that $R(K)$ with $K=2500(\Delta+1)$ is multivariate
zero-free for claw-free graphs of maximum degree $\Delta$, and $R(K)$ with $K=2500\,r(3,k)$ for claw-free graphs of
clique number $k$, where $r(3,k)$ is a Ramsey number. Both parameters are optimal up to a constant factor. For the sector question the answer is negative
when $t\ge2$ and $\Delta\ge6$: for every large $m$ there is a graph of maximum degree $6$ with no induced $S_{1,2,2}$
and a zero $z_m$ with $|2^{-m}z_m-e^{\mathrm i\pi/m}|<m^{-2}$.
\end{abstract}

\noindent\textbf{Mathematics Subject Classification (2020):} 05C31, 05C69, 82B20.

\noindent\textbf{Keywords:} independence polynomial, hard-core model, claw-free graph, zero-free region,
subdivided claw.

\section{Introduction}

For a finite simple graph $G=(V,E)$ and complex weights $\boldsymbol z=(z_v)_{v\in V}$ put
\[
Z_G(\boldsymbol z)=\sum_{I}\prod_{v\in I}z_v ,
\]
the sum running over the independent sets $I$ of $G$, the empty set included. When all weights equal $z$ we write
$Z_G(z)$. A set $F\subseteq\mathbb C$ is a multivariate zero-free region for a class of graphs if
$Z_G(\boldsymbol z)\ne0$ for every $G$ in the class and every $\boldsymbol z\in F^{V(G)}$; the weights are chosen
independently in $F$.

By the Heilmann--Lieb theorem \cite{HL} and its extension by Chudnovsky and Seymour \cite{CS}, the univariate
polynomial $Z_G(z)$ of a claw-free graph has only negative real zeros. Jerrum and Patel \cite{JP} study how far this
extends to graphs with no induced copy of the subdivided claw $S_{i,j,k}$, the tree obtained from the claw by
subdividing its three edges $i-1$, $j-1$ and $k-1$ times. They prove that for fixed $t$ and $\Delta$ the graphs of
maximum degree $\Delta$ with no induced $S_{t,t,t}$ have an open univariate zero-free region containing
$[0,\infty)$ \cite[Theorem~3]{JP}, and an open multivariate zero-free region containing any given interval
$[0,\lambda_0]$ \cite[Theorem~5]{JP}. The second statement cannot be extended to all of $[0,\infty)$: the blown-up
cycles $C_{2n}[2K_1]$, obtained from the cycle $C_{2n}$ by replacing every vertex with two nonadjacent twins, have maximum degree $4$ and no induced $S_{1,2,2}$, and have no open multivariate zero-free
region containing $[0,\infty)$ \cite[Lemma~15]{JP}. As in \cite{JP}, for $K>0$ let
\[
R(K)=\Bigl\{z\in\mathbb C:\ (\Im z)^2<\frac{\Re z}{K}+\frac1{4K^2}\Bigr\}.
\]
By a theorem implicit in the work of Lebowitz, Ruelle and Speer \cite{LRS}, made explicit by Jerrum and Patel,
$R(k^2/4)$ is a multivariate zero-free region for the line graphs of clique number at most $k$ \cite[Theorem~4]{JP}. They prove that $R(k)$ is
a multivariate zero-free region for the claw-free graphs of clique number at most $k$ in which every component has a
simplicial clique \cite[Theorem~6]{JP}. Here a nonempty clique $Q$ is simplicial if for every $u\in Q$ the neighbours
of $u$ outside $Q$ form a clique. Leake and Ryder \cite{LR} proved that the independence polynomial of every claw-free graph is
same-phase stable, that is, for all positive reals $\lambda_v$ the univariate polynomial $z\mapsto Z_G((\lambda_vz)_v)$
has only real zeros. This concerns weights of a common phase and gives no open multivariate region. For claw-free
graphs with a simplicial clique and clique number $\omega$ they show that every zero is at most $-1/(4(\omega-1))$
\cite[Corollary~4.18]{LR}, and the Schl\"afli graph, with $\omega=6$, shows that this bound fails for claw-free graphs
in general \cite[Section~5.1]{LR}. Fialho and Procacci \cite{FP} proved a zero-free disc around $0$ for claw-free
graphs whose radius exceeds Shearer's.

In \cite[Section~7]{JP} Jerrum and Patel ask two questions. The first is whether some form of their Theorem~6 extends
to all claw-free graphs of bounded degree or bounded clique number, and, more basically, whether these classes have an
open multivariate zero-free region containing $[0,\infty)$. The second is whether the region in their Theorem~3 can be
taken to be a sector $\{|\arg z|<\theta\}$ with $\theta>0$ depending on $\Delta$ and $t$. For $S_{1,1,2}$-free graphs
of maximum degree $\Delta$, Bencs \cite{Bencs} proved that $|\arg z|\le\pi/\Delta$ is zero-free (see
\cite[Section~1]{JP}), so the second question has a positive answer for $S_{1,1,2}$-free graphs, in particular for
$t=1$.

The answer to the first question is positive. For claw-free graphs in which every component has a simplicial clique,
Jerrum and Patel prove \cite[Theorem~6]{JP} with the recurrence
\[
W(H,Q)=1+\sum_{u\in Q}\frac{z_u}{W(H-Q,Q_u)},\qquad W(H,Q)=\frac{Z_H}{Z_{H-Q}},
\]
where $Q$ is a simplicial clique of $H$ and $Q_u=N_H(u)\setminus Q$ is empty or a simplicial clique of $H-Q$; by
\cite[Lemma~12]{JP} the ratios $W(H,Q)$ lie in the half-plane $\Re w\ge1/2$ when the clique number is at most $k$
and all weights lie in $R(k)$. They note that deleting a vertex and its
neighbours from a connected claw-free graph leaves a graph in the class of their Theorem~6 \cite[Section~1]{JP}, but
this does not give the answer by itself. If $v$ is a vertex of a claw-free graph $H$ and $B=H-N_H[v]$, then for each
neighbour $i$ of $v$ the set $Q_i=N_H(i)\cap V(B)$ is empty or a simplicial clique of $B$, and
\[
\begin{gathered}
\frac{Z_H}{Z_B}=1+z_v+\sum_{i}z_iA_i+\sum_{\{i,j\}}z_iz_jA_{ij},\\
A_i=\frac1{W(B,Q_i)},\qquad A_{ij}=\frac{A_i}{W(B-Q_i,Q_j\setminus Q_i)}=\frac{A_j}{W(B-Q_j,Q_i\setminus Q_j)},
\end{gathered}
\]
the last sum over the nonadjacent pairs of neighbours of $v$. The half-plane bounds the real parts of the $A_i$ and
$A_{ij}$ but not their arguments, so it gives no lower bound on $\Re(z_iz_jA_{ij})$, and the terms can cancel. We run
the recurrence in the smaller set
\[
\D=\{w\in\mathbb C:\ \Re w\ge1/2+4|\Im w|\},
\]
whose points have modulus at least $1/2$ and argument at most $\arctan(1/4)$ in absolute value, and show that $\D$ is
preserved when every weight lies in $R(2500(\Delta+1))$, that is, when every weight is $z=(\sigma+\mathrm i\tau)^2$
with $\sigma\ge0$ and $|\tau|<\eta=(10000(\Delta+1))^{-1/2}$. Call such a weight large if $\sigma\ge10\eta$ and small
otherwise. A product of large weights in the sum has argument less than $4\arctan(1/4)<\pi/2$, so its real part is
nonnegative. A small weight has modulus at most $101\eta^2$. When $z_i$ is large and $z_j$ small, the first
factorization of $A_{ij}$ gives $|z_iz_jA_{ij}|\le202\eta^2|z_iA_i|$. This is a small multiple of the term $z_iA_i$,
whose real part is at least $\tfrac{15}{17}|z_iA_i|$. The second factorization does the same with $i$ and $j$ exchanged.
Summing,
\[
\Re\frac{Z_H}{Z_B}>0.97
\]
for every vertex $v$ of every induced subgraph $H$, and $Z_G\ne0$ follows by induction.

Here $r(3,k)$ denotes the least $n$ such that every graph on $n$ vertices contains a triangle or an independent set of
size $k$.

\begin{theorem}\label{thm:main}
Let $\Delta\ge0$ and let $G$ be a claw-free graph of maximum degree at most $\Delta$. If $z_v\in R(2500(\Delta+1))$
for every $v\in V(G)$, then $Z_G(\boldsymbol z)\ne0$. Consequently, if $G$ is claw-free of clique number at most
$k\ge1$ and $z_v\in R(2500\,r(3,k))$ for every $v$, then $Z_G(\boldsymbol z)\ne0$.
\end{theorem}

The second statement follows from the first. In a claw-free graph of clique number at most $k$, the neighbourhood
of a vertex induces a graph whose complement is triangle-free and has no independent set of size $k$, so the
neighbourhood has at most $r(3,k)-1$ vertices.
Thus $\Delta+1\le r(3,k)$, and $R(K)\subseteq R(K')$ for $K\ge K'$.

The parameters in \Cref{thm:main} cannot be improved by more than a constant factor.

\begin{proposition}\label{prop:lower}
\begin{enumerate}[label=(\alph*)]
\item If $R(K)$ is a zero-free region for the univariate independence polynomials of the claw-free graphs of maximum
degree at most $\Delta$, then $K\ge(\Delta+1)/4$.
\item If $R(K)$ is a zero-free region for the univariate independence polynomials of the claw-free graphs of clique
number at most $k$, then $K\ge(r(3,k+1)-1)/8$.
\end{enumerate}
\end{proposition}

Adding an isolated vertex to a triangle-free graph on $r(3,k)-1$ vertices with no independent set of size $k$ shows
that $r(3,k+1)-1\ge r(3,k)$. So the parameter in \Cref{prop:lower}(b) is at least $r(3,k)/8$, and the bounds in
\Cref{thm:main} and \Cref{prop:lower}(b) agree up to the factor $20000$. By the theorems of Ajtai, Koml\'os and Szemer\'edi \cite{AKS} and of Kim \cite{Kim}, $r(3,k)$ lies
between two positive constant multiples of $k^2/\log k$. So \Cref{prop:lower}(b) shows that for each fixed $C>0$ the
region $R(Ck)$ is not zero-free for all claw-free graphs of clique number at most $k$ once $k$ is large, even for
univariate weights. The class of \cite[Theorem~6]{JP} behaves differently. It contains the graphs $C_{2n}[K_s]$ of
\cite[Section~5.2]{JP}, of clique number $k=2s$, and for these $R(k/2)$ is the maximal multivariate zero-free region
symmetric about the real axis \cite[Lemma~14]{JP}. So the order $k$ is correct for that class, while on all claw-free
graphs of clique number $k$ the correct order is $k^2/\log k$.

The answer to the second question is negative.

\begin{theorem}\label{thm:sector}
For every $\theta>0$ there are a graph $G$ of maximum degree at most $6$ with no induced copy of $S_{1,2,2}$ and a
complex number $z$ with $0<\arg z<\theta$ and $Z_G(z)=0$. More precisely, there are an absolute constant $m_0$ and,
for every $m\ge m_0$, such a graph $G_m$ on $m(5m+1)$ vertices with a zero $z_m$ satisfying
\[
\bigl|2^{-m}z_m-e^{\mathrm i\pi/m}\bigr|<m^{-2}.
\]
\end{theorem}

Since $S_{1,2,2}$ is an induced subgraph of $S_{t,t,t}$ for every $t\ge2$, the graphs $G_m$ have no induced
$S_{t,t,t}$, and no sector $\{|\arg z|<\theta\}$ is zero-free for the graphs of maximum degree $\Delta\ge6$ with no
induced $S_{t,t,t}$. The zeros $z_m$ tend to infinity, so this does not contradict \cite[Theorem~3]{JP}, and no set
$\{|\arg z|<\theta\}\cup\{|z|<r\}$ with $\theta,r>0$ is zero-free for this class. All vertices of $G_m$ carry the same weight, whereas the obstruction in
\cite[Lemma~15]{JP} uses several different weights. For $3\le\Delta\le5$ and $t\ge2$ the question remains open.

The graph $G_m$ is a blow-up of a cycle of length $2m^2$ whose bags are $K_2\sqcup K_1$, $2K_1$ and $3K_1$. By a
transfer-matrix computation, $Z_{G_m}(2^mw)$ is a nonzero multiple of the trace $\tr(Y^m)$ for a $2\times2$ matrix $Y$ with
$\tr Y=1+w+O(m2^{-m})$ and $\det Y=w+O(m2^{-m})$. Where $(\tr Y)^2=4\cos^2(\pi/(2m))\det Y$ the eigenvalues of $Y$
have ratio $e^{\pm\mathrm i\pi/m}$, so $\tr(Y^m)=0$, and this happens at a point $w_m$ within $m^{-2}$ of
$e^{\mathrm i\pi/m}$.

\section{Proof of \texorpdfstring{\Cref{thm:main}}{Theorem 1.1}}\label{sec:main}

Fix $\Delta\ge0$ and put $L=\Delta+1$ and $\alpha=\arctan(1/4)$, so that $\eta=(10000L)^{-1/2}$ is the number of the
introduction. Every induced subgraph of
$G$ is claw-free with maximum degree at most $\Delta$, so every clique of every induced subgraph has at most $L$
vertices.

\begin{lemma}\label{lem:sqrt}
A complex number $z$ lies in $R(2500L)$ if and only if
\begin{equation}\label{eq:sqrt}
z=(\sigma+\mathrm i\tau)^2\qquad\text{with }\sigma\ge0\text{ and }|\tau|<\eta .
\end{equation}
\end{lemma}

\begin{proof}
Write $z$ as in \eqref{eq:sqrt} with $\sigma\ge0$, so $\Re z=\sigma^2-\tau^2$ and $\Im z=2\sigma\tau$; the number
$\tau^2$ does not depend on which such square root is taken. Since $1/(2500L)=4\eta^2$ and
$1/(4(2500L)^2)=4\eta^4$, the condition $z\in R(2500L)$ reads
$4\sigma^2\tau^2<4\eta^2(\sigma^2-\tau^2)+4\eta^4$, that is, $(\tau^2-\eta^2)(\sigma^2+\eta^2)<0$, which is
$\tau^2<\eta^2$.
\end{proof}

Equivalently, the boundary of $R(K)$ is the image of the line $\Im\zeta=1/(2\sqrt K)$ under $\zeta\mapsto\zeta^2$.

Recall from the introduction the set $\D=\{w\in\mathbb C:\ \Re w\ge1/2+4|\Im w|\}$. Its points satisfy $|w|\ge1/2$ and
$|\arg w|\le\alpha$. For $w\in\mathbb C$ let $\ell_\pm(w)=\Re w\pm4\Im w$; then $w\in\D$ if and only if
$\ell_+(w)\ge1/2$ and $\ell_-(w)\ge1/2$. For $w\ne0$ we have $\ell_\pm(1/w)=\ell_\mp(w)/|w|^2$, and
$|\ell_\pm(c)|\le\sqrt{17}\,|c|$ for every $c$.

\begin{lemma}\label{lem:D}
Let $I$ be a finite set with $|I|\le L$. If $w_j\in\D$ and $z_j\in R(2500L)$ for every $j\in I$, then
$1+\sum_{j\in I}z_j/w_j\in\D$.
\end{lemma}

\begin{proof}
Let $z$ be as in \eqref{eq:sqrt} and $w\in\D$. Then $z=\sigma^2+c$ with
$|c|=|2\mathrm i\sigma\tau-\tau^2|\le2\sigma\eta+\eta^2$, and
\[
\ell_\pm\Bigl(\frac zw\Bigr)=\sigma^2\,\frac{\ell_\mp(w)}{|w|^2}+\ell_\pm\Bigl(\frac cw\Bigr)
\ge\frac{\sigma^2}{2|w|^2}-\sqrt{17}\,\frac{2\sigma\eta+\eta^2}{|w|}.
\]
With $y=\sigma/|w|$ the terms involving $\sigma$ give $y^2/2-2\sqrt{17}\,\eta y\ge-34\eta^2$, and the remaining term
is at least $-2\sqrt{17}\,\eta^2$ because $|w|\ge1/2$. Hence $\ell_\pm(z/w)>-43\eta^2$, and both linear forms of
$1+\sum_j z_j/w_j$ exceed $1-43L\eta^2=1-43/10000>1/2$.
\end{proof}

For an induced subgraph $H$ of $G$ we write $Z_H$ for $Z_H(\boldsymbol z)$ with the inherited weights. For a clique
$Q$ of $H$ with $Z_{H-Q}\ne0$ we use the ratio $W(H,Q)=Z_H/Z_{H-Q}$ of the introduction, which was defined there for
simplicial cliques; thus $W(H,\varnothing)=1$. We prove by induction on $|V(H)|$
the following two statements for every induced subgraph $H$ of $G$:
\begin{enumerate}[label=(\roman*)]
\item $Z_H\ne0$;
\item $W(H,Q)\in\D$ for every simplicial clique $Q$ of $H$ and for $Q=\varnothing$.
\end{enumerate}
For $Q=\varnothing$ statement (ii) holds because $\ell_\pm(1)=1$.
Both hold when $H$ is empty. Let $H$ be nonempty and assume both statements for all smaller induced subgraphs.

\medskip\noindent\emph{Statement (ii).}
Let $Q$ be a nonempty simplicial clique of $H$ and for $u\in Q$ recall $Q_u=N_H(u)\setminus Q$. It is a clique by the
definition of a simplicial clique. It is empty or a simplicial clique of $H-Q$: if $x\in Q_u$ had two nonadjacent neighbours
$a,b$ in $H-Q$ outside $Q_u$, then $a$ and $b$ would not be adjacent to $u$, and $\{x,u,a,b\}$ would induce a claw.
An independent set of $H$ contains at most one vertex of $Q$; if it contains $u\in Q$, its other vertices form an
independent set of $H-Q-Q_u$. Hence
\[
W(H,Q)=1+\sum_{u\in Q}z_u\,\frac{Z_{H-Q-Q_u}}{Z_{H-Q}}=1+\sum_{u\in Q}\frac{z_u}{W(H-Q,Q_u)},
\]
where $Z_{H-Q}\ne0$ and $W(H-Q,Q_u)\in\D$ by the induction hypothesis. This is the recurrence of \cite[Lemma~12]{JP},
with $\D$ in place of the half-plane $\Re w\ge1/2$. Since $|Q|\le L$, \Cref{lem:D} gives $W(H,Q)\in\D$.

\medskip\noindent\emph{Statement (i).}
Fix $v\in V(H)$ and let $U=N_H(v)$, $d=|U|\le\Delta$ and $B=H-N_H[v]$. For $i\in U$ recall $Q_i=N_H(i)\cap V(B)$. Each
$Q_i$ is empty or a simplicial clique of $B$. Indeed, two nonadjacent vertices of $Q_i$ would form, with $v$, three pairwise
nonadjacent neighbours of $i$; and if $x\in Q_i$ had nonadjacent neighbours $a,b$ in $B-Q_i$, then $i,a,b$ would be
three pairwise nonadjacent neighbours of $x$. For $i,j\in U$ the set $Q_j\setminus Q_i$ is empty or a simplicial clique of
$B-Q_i$, because the neighbours of a vertex $x\in Q_j\setminus Q_i$ in $B-Q_i-Q_j$ are among its neighbours in $B$
outside $Q_j$. By the induction hypothesis the independence polynomials of $B$, $B-Q_i$ and $B-(Q_i\cup Q_j)$ are
nonzero, and the numbers $A_i$ and $A_{ij}$ of the introduction are
\begin{equation}\label{eq:A}
\begin{gathered}
A_i=\frac{Z_{B-Q_i}}{Z_B}=\frac1{W(B,Q_i)},\\
A_{ij}=\frac{Z_{B-(Q_i\cup Q_j)}}{Z_B}=\frac{A_i}{W(B-Q_i,Q_j\setminus Q_i)}=\frac{A_j}{W(B-Q_j,Q_i\setminus Q_j)}.
\end{gathered}
\end{equation}
Since every element of $\D$ has modulus at least $1/2$ and argument at most $\alpha$ in absolute value,
\begin{equation}\label{eq:Abounds}
|A_i|\le2,\quad|\arg A_i|\le\alpha,\qquad|A_{ij}|\le2\min(|A_i|,|A_j|),\quad|\arg A_{ij}|\le2\alpha .
\end{equation}
Because $H$ is claw-free, $U$ contains no three pairwise nonadjacent vertices. Classifying the independent sets of
$H$ by their intersection with $N_H[v]$, which is empty, $\{v\}$, $\{i\}$ with $i\in U$, or $\{i,j\}$ with
$i,j\in U$ nonadjacent, we obtain
\begin{equation}\label{eq:root}
\frac{Z_H}{Z_B}=1+z_v+\sum_{i\in U}z_iA_i+\sum_{\{i,j\}}z_iz_jA_{ij},
\end{equation}
the last sum over unordered pairs of nonadjacent vertices of $U$.

Write each weight $z$ as in \eqref{eq:sqrt}. Recall that $z$ is large if $\sigma\ge10\eta$ and small
otherwise. For a large weight put $z^+=z$ and $z^-=0$; for a small one put $z^+=0$ and $z^-=z$. If $z$ is large then
$|\arg(\sigma+\mathrm i\tau)|\le\arctan(1/10)$, so $|\arg z^+|\le2\arctan(1/10)<\alpha$; if $z$ is small then
$|z^-|=\sigma^2+\tau^2\le101\eta^2=:\varepsilon_0$. Do this for $z_v$ and for each $z_i$, and put $b_i=|z_i^+A_i|$.

By \eqref{eq:Abounds}, a nonzero term $z_i^+A_i$ has argument at most $2\alpha$ in absolute value, so
$\Re(z_i^+A_i)\ge\cos(2\alpha)\,b_i=\tfrac{15}{17}\,b_i$. A nonzero term $z_i^+z_j^+A_{ij}$ has argument at most
$4\alpha<\pi/2$ in absolute value, so its real part is nonnegative, and so is $\Re z_v^+$. The remaining terms are
bounded using \eqref{eq:A} and \eqref{eq:Abounds}: $|z_v^-|\le\varepsilon_0$, $|z_i^-A_i|\le2\varepsilon_0$,
$|z_i^+z_j^-A_{ij}|\le2\varepsilon_0\,b_i$ by the first factorization of $A_{ij}$, $|z_i^-z_j^+A_{ij}|\le2\varepsilon_0\,b_j$
by the second, and $|z_i^-z_j^-A_{ij}|\le4\varepsilon_0^2$. Each $i\in U$ lies in at most $d-1$ pairs, so taking real
parts in \eqref{eq:root} gives
\[
\Re\frac{Z_H}{Z_B}\ \ge\ 1-(1+2d)\varepsilon_0-4\binom d2\varepsilon_0^2+\sum_{i\in U}\Bigl(\frac{15}{17}-2d\varepsilon_0\Bigr)b_i .
\]
As $d+1\le L$ we have $(1+2d)\varepsilon_0<2L\cdot101\eta^2=0.0202$, $4\binom d2\varepsilon_0^2<0.0202^2/2$ and
$2d\varepsilon_0<0.0202$. Hence $\Re(Z_H/Z_B)>0.97$, and $Z_H\ne0$ because $Z_B\ne0$. This completes the induction,
and statement (i) for $H=G$ is \Cref{thm:main}.

\begin{remark}
On a graph containing an induced claw, \eqref{eq:root} acquires terms with three neighbours of $v$ and the argument
fails, as it must by \cite[Lemma~15]{JP}.
\end{remark}

\section{Proof of \texorpdfstring{\Cref{prop:lower}}{Proposition 1.2}}\label{sec:lower}

For $K>0$ we have $R(K)\cap\mathbb R=(-1/(4K),\infty)$.

(a) The complete graph $K_{\Delta+1}$ is claw-free with maximum degree $\Delta$, and
$Z_{K_{\Delta+1}}(z)=1+(\Delta+1)z$ vanishes at $z=-1/(\Delta+1)$. If $R(K)$ is zero-free then
$-1/(\Delta+1)\le-1/(4K)$, that is, $K\ge(\Delta+1)/4$.

(b) Let $N=r(3,k+1)-1$. By the definition of $r(3,k+1)$ there is a triangle-free graph $\Gamma$ on $N$ vertices with no
independent set of size $k+1$. Its complement $\overline \Gamma$ has clique number at most $k$ and no independent set of
size $3$, so it is claw-free. The independent sets of $\overline \Gamma$ are the cliques of $\Gamma$, namely the empty set, the
$N$ vertices and the $e(\Gamma)$ edges of $\Gamma$, so $Z_{\overline \Gamma}(z)=1+Nz+e(\Gamma)z^2$. For every edge $xy$ of $\Gamma$ the
neighbourhoods of $x$ and $y$ are disjoint, so $\deg x+\deg y\le N$; summing over edges gives
$\sum_x(\deg x)^2\le Ne(\Gamma)$, and by Cauchy--Schwarz $4e(\Gamma)^2=(\sum_x\deg x)^2\le N\sum_x(\deg x)^2\le N^2e(\Gamma)$, so
$4e(\Gamma)\le N^2$. If $e(\Gamma)>0$ the zeros of $1+Nz+e(\Gamma)z^2$ are real and negative and the one nearer to $0$ is
$-2/(N+\sqrt{N^2-4e(\Gamma)})\ge-2/N$; if $e(\Gamma)=0$ the only zero is $-1/N$. If $R(K)$ is zero-free then $-2/N\le-1/(4K)$,
that is, $K\ge N/8$.

\section{Proof of \texorpdfstring{\Cref{thm:sector}}{Theorem 1.3}}\label{sec:sector}

\subsection{The graphs}

Let $m\ge2$ and $h\ge1$. Take a cycle with $2mh$ sites, numbered cyclically, and replace each site by a graph, its
bag, as follows. Odd sites receive $K_2\sqcup K_1$ (an edge and an isolated vertex). Among the even sites, every
$m$-th one receives $3K_1$ and the others receive $2K_1$, so that each run of $2m$ consecutive sites contains $m$
copies of $K_2\sqcup K_1$, $m-1$ copies of $2K_1$ and one copy of $3K_1$. Every vertex of a bag is joined to every
vertex of the two neighbouring bags, and there are no other edges between bags. Let $G_{m,h}$ be the resulting graph;
it has $h(5m+1)$ vertices.

A vertex of an even bag has the $3+3$ vertices of the two neighbouring odd bags as its neighbours. An odd bag has at
most one neighbouring bag of size $3$ because $m\ge2$, so its vertices have degree at most $1+2+3=6$. Thus $G_{m,h}$
has maximum degree at most $6$.

\begin{lemma}\label{lem:E}
$G_{m,h}$ has no induced copy of $S_{1,2,2}$.
\end{lemma}

\begin{proof}
Every bag $X$ is a module: each vertex outside $X$ is adjacent to all of $X$ or to none of it. Let $T$ be an induced
subgraph of $G_{m,h}$ isomorphic to $S_{1,2,2}$. Then $V(T)\cap X$ is a module of $T$, and it is proper since $T$ has
six vertices and $X$ at most three. We claim that a proper module $J$ of a tree $T$ with $|J|\ge2$ consists of
leaves with a common neighbour. As $T$ is connected, some vertex $p\notin J$ has a neighbour in $J$, hence is adjacent
to all of $J$. An edge inside $J$ would form a triangle with $p$, and a second vertex outside $J$ adjacent to $J$ would
form a $4$-cycle with $p$ and two vertices of $J$. So $J$ is independent and $p$ is the only neighbour of each of its
vertices. In $S_{1,2,2}$ the three leaves have three different neighbours, so each bag contains at most one vertex of
$T$. Two vertices in different bags are adjacent exactly when their bags are neighbouring sites of the cycle, so $T$
would be isomorphic to an induced subgraph of a cycle and have maximum degree at most $2$, while $S_{1,2,2}$ has a
vertex of degree $3$.
\end{proof}

\subsection{Transfer matrices}

An independent set of $G_{m,h}$ meets an independent set of sites of the cycle, and inside each such site it is a
nonempty independent set of the bag. Summing over the latter, as in \cite[Section~5.2]{JP}, gives
$Z_{G_{m,h}}(z)=Z_{C_{2mh}}(\boldsymbol x)$, where the cycle carries at each site the weight
\[
a=2z^2+3z,\qquad b=z^2+2z,\qquad c=z^3+3z^2+3z
\]
according as the bag is $K_2\sqcup K_1$, $2K_1$ or $3K_1$. For $x\in\mathbb C$ let
$M(x)=\begin{pmatrix}1&x\\1&0\end{pmatrix}$, with rows indexed by the state of the previous site and columns by the
state of the current one ($0$ empty, $1$ occupied). As in \cite[Section~5.1]{JP}, the independence polynomial of a
weighted cycle is the trace of the product of these matrices taken around the cycle. We have
$M(a)M(b)=\begin{pmatrix}1+a&b\\1&b\end{pmatrix}$ and $M(a)M(c)=\begin{pmatrix}1+a&c\\1&c\end{pmatrix}$.

Let $z\notin\{0,-2\}$ and put
\[
\rho=\frac ab=2-\frac1{z+2},\qquad s=\frac cb=z+1+\frac1{z+2},\qquad\varepsilon=\frac1b,\qquad
P=\begin{pmatrix}\rho+\varepsilon&1\\\varepsilon&1\end{pmatrix}.
\]
Then $M(a)M(b)=bP$ and $M(a)M(c)=bP\,\mathrm{diag}(1,s)$, so with
\[
Y=2^{-m}P^m\,\mathrm{diag}(1,s)
\]
we get
\begin{equation}\label{eq:trace}
Z_{G_{m,h}}(z)=\bigl(2^mb^m\bigr)^h\tr(Y^h).
\end{equation}

\begin{lemma}\label{lem:asym}
There is an absolute constant $m_0$ such that for $m\ge m_0$, uniformly for $1/2\le|w|\le2$ and $z=2^mw$,
\[
\tr Y=1+w+O(m2^{-m}),\qquad\det Y=w+O(m2^{-m}).
\]
\end{lemma}

\begin{proof}
All constants below are absolute, and $m$ is large. We have $|z|\ge2^{m-1}$, so $1/(z+2)=O(2^{-m})$,
$\varepsilon=O(4^{-m})$, $s=O(2^m)$, $|\rho|\ge3/2$, $|\rho-1|\ge1/2$ and $(\rho/2)^m=(1-1/(2(z+2)))^m=1+O(m2^{-m})$; in
particular $|\rho|^m/2^m=O(1)$.

Since $\det P=\rho$, we get $\det Y=2^{-2m}\rho^ms=(\rho/2)^m\,s2^{-m}$, and $s2^{-m}=w+O(2^{-m})$; this gives the second
estimate.

For the trace let $P_0=\begin{pmatrix}\rho&1\\0&1\end{pmatrix}$, so that $P-P_0$ has maximum absolute row sum
$\|P-P_0\|_\infty=|\varepsilon|$ and
\[
P_0^j=\begin{pmatrix}\rho^j&(\rho^j-1)/(\rho-1)\\0&1\end{pmatrix},\qquad\|P_0^j\|_\infty\le5|\rho|^j\quad(j\ge0).
\]
Expanding $P^m=(P_0+(P-P_0))^m$, a product containing $n\ge1$ factors $P-P_0$ has norm at most
$5^{n+1}|\rho|^{m-n}|\varepsilon|^n$, and there are $\binom mn$ such products, so
\[
\|P^m-P_0^m\|_\infty\le5|\rho|^m\Bigl[\Bigl(1+\frac{5|\varepsilon|}{|\rho|}\Bigr)^m-1\Bigr]=O(m|\varepsilon|\,|\rho|^m),
\]
the last step because $m|\varepsilon|=O(m4^{-m})$ is small. Multiplying by $\mathrm{diag}(1,s)$, whose norm is
$O(2^m)$, taking the trace and dividing by $2^m$ therefore changes the result by
$O(m|\varepsilon|\,|\rho|^m)=O(m4^{-m}2^m)=O(m2^{-m})$. Finally
$2^{-m}\tr\bigl(P_0^m\,\mathrm{diag}(1,s)\bigr)=(\rho/2)^m+s2^{-m}=1+w+O(m2^{-m})$.
\end{proof}

\subsection{The zeros}

Take $h=m$ and let $G_m=G_{m,m}$, a graph on $m(5m+1)$ vertices. On the annulus $1/2\le|w|\le2$ the entries of $Y$
are holomorphic functions of $w$, and so is
\[
\Phi_m(w)=(\tr Y)^2-4\cos^2\Bigl(\frac\pi{2m}\Bigr)\det Y .
\]
By \Cref{lem:asym}, $\Phi_m$ differs by $O(m2^{-m})$ from
\[
(1+w)^2-4\cos^2\Bigl(\frac\pi{2m}\Bigr)w=(w-e^{\mathrm i\pi/m})(w-e^{-\mathrm i\pi/m}).
\]
On the circle $|w-e^{\mathrm i\pi/m}|=m^{-2}$, which lies in the annulus, the right-hand side has modulus at least
$m^{-2}(2\sin(\pi/m)-m^{-2})\ge m^{-3}$ for large $m$, while $m2^{-m}=o(m^{-3})$. By Rouch\'e's theorem $\Phi_m$ has a
zero $w_m$ with $|w_m-e^{\mathrm i\pi/m}|<m^{-2}$.

At $w_m$ we have $\det Y=w_m+O(m2^{-m})\ne0$, so the eigenvalues $\lambda_1,\lambda_2$ of $Y$ are nonzero. Their ratio
$q=\lambda_1/\lambda_2$ satisfies $(1+q)^2/q=(\tr Y)^2/\det Y=4\cos^2(\pi/(2m))$, that is,
$q+q^{-1}=2\cos(\pi/m)$, so $q=e^{\pm\mathrm i\pi/m}$. In particular $\lambda_1\ne\lambda_2$, $Y$ is diagonalizable, and
$\lambda_1^m=-\lambda_2^m$, so $\tr(Y^m)=\lambda_1^m+\lambda_2^m=0$. The prefactor in \eqref{eq:trace} does not
vanish at $z_m=2^mw_m$ because $b=z_m(z_m+2)\ne0$, so $Z_{G_m}(z_m)=0$. Finally
$\arg z_m=\arg w_m=\pi/m+O(m^{-2})$ and $|z_m|2^{-m}\to1$. Together with \Cref{lem:E} and the degree bound this proves
\Cref{thm:sector}.

\begin{remark}
The error term is small already for moderate $m$. To six decimals, the zero of $Z_{G_m}$ near $2^me^{\mathrm i\pi/m}$ has
argument $0.196525$ for $m=16$ and $0.157093$ for $m=20$, against $\pi/16\approx0.196350$ and
$\pi/20\approx0.157080$.
\end{remark}

\section*{Acknowledgments}

This work was done by multiple AI agents, Claude (Anthropic) and Codex (OpenAI), connected to Math
Brain~\cite{MathBrainRepo} and working together inside the Brain framework~\cite{BrainFramework}, a shared, persistent
memory in which the agents record claims, proofs, failed attempts and methods, retrieve one another's results,
coordinate who works on what, and audit each other's work. For more on the Brain and Math Brain, see the Math Brain
paper~\cite{MathBrainPaper}, the Math Brain repository~\cite{MathBrainRepo} and the Brain framework
repository~\cite{BrainFramework}.

\section*{Funding}

The research presented in this paper was supported by the European Research Council (ERC) under the European
Union's Horizon Europe research and innovation programme (grant agreement No.\ 101041711). It was also supported by the Simons
Foundation (as part of the Collaboration on the Mathematical and Scientific Foundations of Deep Learning), by Heights Labs, by
Convex Nexus Capital, and by the Israel Science Foundation (grants 2258/19 and 4101/25).

\section*{Declaration of competing interest}

The author holds equity in Heights Labs and Convex Nexus Capital and has a management role at Convex Nexus Capital.

\section*{Declaration of generative AI and AI-assisted technologies}

The proofs of \Cref{thm:main}, \Cref{prop:lower} and \Cref{thm:sector} were developed with AI systems (OpenAI Codex
and Anthropic Claude agents working under the author's direction), which also drafted the text and wrote the computer
checks. The questions are those of Jerrum and Patel \cite[Section~7]{JP}. The regions $R(K)$, the notion of a
simplicial clique and the recurrence for $W(H,Q)$ in \Cref{sec:main} are taken from \cite{JP} (Theorem~6 and the proof
of Lemma~12 there), and the transfer matrices and the summation over bags in \Cref{sec:sector} follow
\cite[Sections~5.1 and~5.2]{JP}. The region $R(k^2/4)$ for line graphs is due to Lebowitz, Ruelle and Speer
\cite{LRS}, the Schl\"afli graph example and same-phase stability to Leake and Ryder \cite{LR}, and the order of
magnitude of $r(3,k)$ to Ajtai, Koml\'os and Szemer\'edi \cite{AKS} and Kim \cite{Kim}. The author directed the work and is responsible for the content of the paper and for
decisions concerning its dissemination.

\section*{Data and code availability}

The Python script at
\begin{center}
\url{https://github.com/yspennstate/claw-free-independence-zeros/tree/v1.2}
\end{center}
compares the transfer-matrix formula \eqref{eq:trace} with a direct evaluation of the independence polynomial on small
graphs $G_{m,h}$. It checks by exhaustive search that $G_{2,2}$ has no induced $S_{1,2,2}$ and computes the zeros near
$2^me^{\mathrm i\pi/m}$ for $m=8,12,16,20,30$. It also evaluates $\Re(Z_H/Z_B)$ in \eqref{eq:root} on random claw-free graphs
with weights in $R(2500(\Delta+1))$, and computes the independence polynomials of the examples in the proof of
\Cref{prop:lower} from the graphs.

\begin{thebibliography}{99}

\bibitem{AKS}
M.~Ajtai, J.~Koml\'os and E.~Szemer\'edi,
\emph{A note on Ramsey numbers},
J.\ Combin.\ Theory Ser.\ A \textbf{29} (1980), 354--360.
\doi{10.1016/0097-3165(80)90030-8}

\bibitem{Bencs}
F.~Bencs,
\emph{Graphs, groups and measures},
Ph.D.\ thesis, Central European University, 2019.

\bibitem{CS}
M.~Chudnovsky and P.~Seymour,
\emph{The roots of the independence polynomial of a clawfree graph},
J.\ Combin.\ Theory Ser.\ B \textbf{97} (2007), 350--357.
\doi{10.1016/j.jctb.2006.06.001}

\bibitem{FP}
P.~M.~S.~Fialho and A.~Procacci,
\emph{On the independent set polynomial of graphs and claw-free graphs},
preprint, 2025, arXiv:2505.22766.

\bibitem{HL}
O.~J.~Heilmann and E.~H.~Lieb,
\emph{Theory of monomer-dimer systems},
Comm.\ Math.\ Phys.\ \textbf{25} (1972), 190--232.
\doi{10.1007/BF01877590}

\bibitem{JP}
M.~Jerrum and V.~Patel,
\emph{Zero-free regions for the independence polynomial on restricted graph classes},
J.\ London Math.\ Soc.\ (2) \textbf{113} (2026), no.~2, e70458.
\doi{10.1112/jlms.70458}; arXiv:2510.01466v2.

\bibitem{Kim}
J.~H.~Kim,
\emph{The Ramsey number $R(3,t)$ has order of magnitude $t^2/\log t$},
Random Structures Algorithms \textbf{7} (1995), 173--207.
\doi{10.1002/rsa.3240070302}

\bibitem{LR}
J.~D.~Leake and N.~R.~Ryder,
\emph{Generalizations of the matching polynomial to the multivariate independence polynomial},
Algebr.\ Comb.\ \textbf{2} (2019), no.~5, 781--802.
\url{https://www.numdam.org/item/ALCO_2019__2_5_781_0/}

\bibitem{LRS}
J.~L.~Lebowitz, D.~Ruelle and E.~R.~Speer,
\emph{Location of the Lee--Yang zeros and absence of phase transitions in some Ising spin systems},
J.\ Math.\ Phys.\ \textbf{53} (2012), 095211.
\doi{10.1063/1.4738622}

\bibitem{MathBrainRepo}
Y.~Shmalo, \emph{Math Brain: public snapshot, code and tests}, 2026.
\url{https://github.com/yspennstate/math-brain-public}

\bibitem{MathBrainPaper}
Y.~Shmalo, \emph{Math Brain: shared, auditable memory for multi-agent mathematical research}, October 2026.
\url{https://github.com/yspennstate/math-brain-public/tree/main/papers/math-brain-2026-10-24}

\bibitem{BrainFramework}
Y.~Shmalo, \emph{Brain Framework}, public software repository, 2026.
\url{https://github.com/yspennstate/brain-framework-public}

\end{thebibliography}

\end{document}
